\section{InstructorNotes.sty}\label{sec:ref-instructornotes}

The instructor-only content layer. Everything here is gated on one flag,
\cs{ifinstructornotes}, and \emph{fully gobbled} when off -- a student
build carries no trace. The flag defaults on when \cs{ShowInstructorNotes} is
defined, which the build's \env{instructor} view injects
(section~\ref{sec:build-and-find}). One staff tier only: the instructor
copy \emph{is} the staff copy (TAs included); anything unfit for a TA stays
out of the source.

\subsection*{The three shapes}

\begin{center}
\begin{tabularx}{\linewidth}{@{}l Y@{}}
\cs{InstructorInline}\texttt{\{...\}} & slate run-in inside the paragraph;
  wraps across lines. For quick in-flow asides -- anything you will want to
  \emph{find} again belongs in a box or the margin.\\
\env{InstructorBox}\texttt{[marking]} & the side-barred block: pale slate
  fill, a slate bar in the left margin with a rotated word -- \emph{note}
  (default) or \emph{marking}. The word differentiates, not the colour.\\
\cs{InstructorMargin}\texttt{\{...\}} & bold slate sans in the margin
  (via \cs{marginnote}, so it never floats).\\
\end{tabularx}
\end{center}

Everything is slate (\env{InstructorInk}) on purpose: staff content is
``editorial, behind the scenes,'' one neutral colour with one meaning --
distinct from the student palette and from green $=$ solutions.

\subsection*{The page frame}

In an instructor build, every page carries a slate frame just inside the
page edge with an \emph{Instructor copy} tab at bottom right (clear of the
centred page-number footers). An ever-present, low-distraction marker --
it replaced both a title stamp (visible only on page 1) and a watermark
(fights the content). Knobs: \cs{InstructorFrameInset} (edge gap, 2.5\,mm)
and \cs{InstructorFrameRule} (thickness, 1.5\,mm).

Note the frame installs \emph{at load time}, keyed to the injected flag --
which is why this manual demonstrates the boxes by flipping
\cs{instructornotestrue} locally (the environments test the flag at use
time) without ever defining \cs{ShowInstructorNotes}.

\subsection*{Interaction with Workspace}

\env{InstructorBox} routes through \cs{WSZeroHeight}: in overlay mode it is
a zero-height overlay like \env{solution} (section~\ref{sec:overlay}). The
package stays self-contained via a fallback -- under CourseDocument, which
has no Workspace, boxes simply render at normal height.
